\documentclass[10pt,a4paper]{amsart}
\usepackage{amsmath,amssymb,amsthm,graphicx}
\usepackage{verbatim}
\newenvironment{reference}{\comment}{\endcomment}
%\newenvironment{reference}{}{}
\newenvironment{slogan}{\comment}{\endcomment}
\newenvironment{history}{\comment}{\endcomment}
% For commutative diagrams we use Xy-pic
\usepackage[all]{xy}

% We use 2cell for 2-commutative diagrams.
\xyoption{2cell}
% Theorem environments.
%
\theoremstyle{plain}
\newtheorem{theorem}[subsection]{Theorem}
\newtheorem{proposition}[subsection]{Proposition}
\newtheorem{lemma}[subsection]{Lemma}

\theoremstyle{definition}
\newtheorem{definition}[subsection]{Definition}
\newtheorem{example}[subsection]{Example}
\newtheorem{exercise}[subsection]{Exercise}
\newtheorem{situation}[subsection]{Situation}

\theoremstyle{remark}
\newtheorem{remark}[subsection]{Remark}
\newtheorem{remarks}[subsection]{Remarks}

\numberwithin{equation}{subsection}

% Macros
%
\def\lim{\mathop{\mathrm{lim}}\nolimits}
\def\colim{\mathop{\mathrm{colim}}\nolimits}
\def\Spec{\mathop{\mathrm{Spec}}}
\def\Hom{\mathop{\mathrm{Hom}}\nolimits}
\def\Ext{\mathop{\mathrm{Ext}}\nolimits}
\def\SheafHom{\mathop{\mathcal{H}\!\mathit{om}}\nolimits}
\def\SheafExt{\mathop{\mathcal{E}\!\mathit{xt}}\nolimits}
\def\Sch{\mathit{Sch}}
\def\Mor{\mathop{\mathrm{Mor}}\nolimits}
\def\Ob{\mathop{\mathrm{Ob}}\nolimits}
\def\Sh{\mathop{\mathit{Sh}}\nolimits}
\def\NL{\mathop{N\!L}\nolimits}
\def\CH{\mathop{\mathrm{CH}}\nolimits}
\def\proetale{{pro\text{-}\acute{e}tale}}
\def\etale{{\acute{e}tale}}
\def\QCoh{\mathit{QCoh}}
\def\Ker{\mathop{\mathrm{Ker}}}
\def\Im{\mathop{\mathrm{Im}}}
\def\Coker{\mathop{\mathrm{Coker}}}
\def\Coim{\mathop{\mathrm{Coim}}}

% Boxtimes
%
\DeclareMathSymbol{\boxtimes}{\mathbin}{AMSa}{"02}

%
% Macros for moduli stacks/spaces
%
\def\QCohstack{\mathcal{QC}\!\mathit{oh}}
\def\Cohstack{\mathcal{C}\!\mathit{oh}}
\def\Spacesstack{\mathcal{S}\!\mathit{paces}}
\def\Quotfunctor{\mathrm{Quot}}
\def\Hilbfunctor{\mathrm{Hilb}}
\def\Curvesstack{\mathcal{C}\!\mathit{urves}}
\def\Polarizedstack{\mathcal{P}\!\mathit{olarized}}
\def\Complexesstack{\mathcal{C}\!\mathit{omplexes}}
% \Pic is the operator that assigns to X its picard group, usage \Pic(X)
% \Picardstack_{X/B} denotes the Picard stack of X over B
% \Picardfunctor_{X/B} denotes the Picard functor of X over B
\def\Pic{\mathop{\mathrm{Pic}}\nolimits}
\def\Picardstack{\mathcal{P}\!\mathit{ic}}
\def\Picardfunctor{\mathrm{Pic}}
\def\Deformationcategory{\mathcal{D}\!\mathit{ef}}
\usepackage{cleveref}
\crefname{originallemma}{original lemma}{original lemmas}
\Crefname{originallemma}{Original lemma}{Original lemmas}
\makeatletter
\newlabel{lemma-functoriality-differentials-ringed-spaces}{{lemma-functoriality-differentials-ringed-spaces}{source}{}{sourcepointer-1}{}}
\newlabel{lemma-functoriality-differentials-ringed-spaces@cref}{{[originallemma][0][]lemma-functoriality-differentials-ringed-spaces}{[1][0][]source}{}{}{}}
\makeatother
\crefname{originallemma}{original lemma}{original lemmas}
\Crefname{originallemma}{Original lemma}{Original lemmas}
\makeatletter
\newlabel{lemma-stalk-module-differentials}{{lemma-stalk-module-differentials}{source}{}{sourcepointer-2}{}}
\newlabel{lemma-stalk-module-differentials@cref}{{[originallemma][0][]lemma-stalk-module-differentials}{[1][0][]source}{}{}{}}
\makeatother
\crefname{originalreference}{original reference}{original references}
\Crefname{originalreference}{Original reference}{Original references}
\makeatletter
\newlabel{algebra-lemma-exact-sequence-differentials}{{algebra-lemma-exact-sequence-differentials}{source}{}{sourcepointer-3}{}}
\newlabel{algebra-lemma-exact-sequence-differentials@cref}{{[originalreference][0][]algebra-lemma-exact-sequence-differentials}{[1][0][]source}{}{}{}}
\makeatother
\crefname{originalreference}{original reference}{original references}
\Crefname{originalreference}{Original reference}{Original references}
\makeatletter
\newlabel{sheaves-lemma-stalk-pullback-modules}{{sheaves-lemma-stalk-pullback-modules}{source}{}{sourcepointer-4}{}}
\newlabel{sheaves-lemma-stalk-pullback-modules@cref}{{[originalreference][0][]sheaves-lemma-stalk-pullback-modules}{[1][0][]source}{}{}{}}
\makeatother
\title{The exact sequence of modules of differentials for a composition of morphisms of ringed spaces}
\author{The Stacks Project Authors}
\date{stacks-project:modules.tex:local-snapshot-20261012; GNU FDL 1.2 or later}
\begin{document}
\maketitle
\noindent\textit{Editorial note.} All delivered mathematical text is the author's own, copied from the pinned Stacks Project chapter source sources/stacks-modules/pinned-original-source.tex; no statement or proof step was rewritten, reordered or removed.\par\medskip
\noindent\textit{Editorial note.} Preamble provenance: the Stacks Project keeps one single shared preamble for every chapter; this card takes it directly from the staged sources/stacks-modules/preamble.tex as three verbatim spans (lines 8-12, 14-18 and 154-217), so no derived file was written and nothing was added or redefined. The wrapper class and the base packages are supplied by the assembler.\par\medskip
\noindent\textit{Editorial note.} Bibliography provenance: neither the counted statement nor its proof contains a citation (this lane rescans the whole span before writing the card), so the card carries no bibliography and no entry was reconstructed.\par\medskip
\noindent\textit{Editorial note.} Selection provenance: this card was opened under the file-based control contract (CONTROL.md), with the mechanical pre-card gate run on its statement and proof interval before the card was written.\par\medskip
\noindent\textit{Allowed standard prerequisites (editorial).}\begin{itemize}
\item The functoriality of the module of differentials for a morphism of ringed spaces, together with the description of the pullback of a module on stalks. Both are invoked inside the counted proof (source lines 5341-5342, 5354-5361) as pointers to a lemma of this chapter and a result of another chapter. Identifies the maps of the sequence with the maps appearing in the algebraic exact sequence.
\item The compatibility of the formation of differentials with taking stalks. Invoked inside the counted proof at source line 5347 as a pointer to a lemma of this chapter. Turns the sequence of sheaves into a sequence of modules over the local rings of the stalks.
\end{itemize}\medskip
\ifdefined\hypertarget\hypertarget{sourcepointer-1}{}\fi
\noindent\textit{Source reference.} Original lemma lemma-functoriality-differentials-ringed-spaces: sources/stacks-modules/pinned-original-source.tex, source lines 5323--5323. Navigation pointer only: the counted text refers to the Stacks item lemma-functoriality-differentials-ringed-spaces, which belongs to another chapter of the project and is not part of this card. The pointer carries the target's own label as its number because this chapter's staged snapshot does not contain the other chapter's tag table; its content is not reproduced here.\par\medskip
\ifdefined\hypertarget\hypertarget{sourcepointer-2}{}\fi
\noindent\textit{Source reference.} Original lemma lemma-stalk-module-differentials: sources/stacks-modules/pinned-original-source.tex, source lines 5197--5197. Navigation pointer only: the counted text refers to the Stacks item lemma-stalk-module-differentials, which belongs to another chapter of the project and is not part of this card. The pointer carries the target's own label as its number because this chapter's staged snapshot does not contain the other chapter's tag table; its content is not reproduced here.\par\medskip
\ifdefined\hypertarget\hypertarget{sourcepointer-3}{}\fi
\noindent\textit{Source reference.} Original reference algebra-lemma-exact-sequence-differentials: sources/stacks-modules/pinned-original-source.tex, source lines 5355--5355. Navigation pointer only: the counted text refers to the Stacks item algebra-lemma-exact-sequence-differentials, which belongs to another chapter of the project and is not part of this card. The pointer carries the target's own label as its number because this chapter's staged snapshot does not contain the other chapter's tag table; its content is not reproduced here.\par\medskip
\ifdefined\hypertarget\hypertarget{sourcepointer-4}{}\fi
\noindent\textit{Source reference.} Original reference sheaves-lemma-stalk-pullback-modules: sources/stacks-modules/pinned-original-source.tex, source lines 1418--1418. Navigation pointer only: the counted text refers to the Stacks item sheaves-lemma-stalk-pullback-modules, which belongs to another chapter of the project and is not part of this card. The pointer carries the target's own label as its number because this chapter's staged snapshot does not contain the other chapter's tag table; its content is not reproduced here.\par\medskip
\begin{lemma}
\label{lemma-triangle-differentials}
Let $f : X \to Y$, $g : Y \to S$ be morphisms of ringed spaces
Then there is a canonical exact sequence
$$
f^*\Omega_{Y/S} \to \Omega_{X/S} \to \Omega_{X/Y} \to 0
$$
where the maps come from applications of
Lemma \ref{lemma-functoriality-differentials-ringed-spaces}.
\end{lemma}

\begin{proof}
By taking induced maps in stalks at $x \in X$ and using
Lemma \ref{lemma-stalk-module-differentials} we obtain a sequence
$$
\mathcal{O}_{X, x} \otimes_{\mathcal{O}_{Y, f(x)}}
\Omega_{\mathcal{O}_{Y, f(x)}/\mathcal{O}_{S, g(f(x))}} \to
\Omega_{\mathcal{O}_{X, x}/\mathcal{O}_{S, g(f(x))}} \to
\Omega_{\mathcal{O}_{X, x}/\mathcal{O}_{Y, f(x)}} \to 0
$$
It suffices to see that the maps of the sequence are the same as
the ones in Algebra, Lemma \ref{algebra-lemma-exact-sequence-differentials}.
This is because the characterization of the maps in
Lemma \ref{lemma-functoriality-differentials-ringed-spaces}
and because via the isomorphism of
Sheaves, Lemma \ref{sheaves-lemma-stalk-pullback-modules}
we have $(f^*s)_x = s_x \otimes 1$,
for any local section $s$ of a sheaf of $\mathcal{O}_Y$-modules.
\end{proof}

\end{document}
